\documentclass[11pt,a4paper]{article}

\usepackage[margin=2cm]{geometry}
\usepackage{enumitem}
\usepackage[hidelinks]{hyperref}
\usepackage[T1]{fontenc}
\usepackage[utf8]{inputenc}
\usepackage{titlesec}

\setlength{\parindent}{0pt}
\pagenumbering{gobble}

\titleformat{\section}{\large\bfseries}{}{0em}{}[\titlerule]

\begin{document}

% ================= HEADER =================
\begin{center}
{\Large \textbf{MIGUEL ANGEL ANAY GOMEZ}}\\
Data Engineer | Cloud Data | ETL | Azure | AWS | Python | SQL\\
\vspace{4pt}
\href{mailto:contacto@miguel-anay.nom.pe}{contacto@miguel-anay.nom.pe} |
\href{https://www.linkedin.com/in/miguel-anay}{LinkedIn} |
Perú |
936 327 402 |
\href{https://portfolio.miguel-anay.nom.pe}{Portfolio}
\end{center}

\vspace{6pt}

% ================= PERFIL =================
\section*{Perfil Profesional}

Ingeniero de Software con más de \textbf{10 años de experiencia}, especializado en \textbf{Data Engineering y arquitecturas Cloud}.  
Experiencia en el diseño, desarrollo y despliegue de \textbf{pipelines de datos}, procesos \textbf{ETL/ELT} y soluciones \textbf{Data Lake} sobre entornos \textbf{Azure y AWS}.  
Sólidos conocimientos en \textbf{Python y SQL}, optimización de performance, control de calidad de datos, monitoreo de consumo en la nube y despliegue de soluciones escalables orientadas a analítica y explotación de información.  
Experiencia práctica utilizando \textbf{Databricks, Apache Spark y CI/CD} en flujos de datos productivos.

\vspace{6pt}

% ================= EXPERIENCIA =================
\section*{Experiencia Profesional}

\textbf{CEO \& Data / Full Stack Engineer – Yiwu Import Corporation} \hfill \textit{Abr 2025 – Actual}
\begin{itemize}[leftmargin=*]
  \item Diseño y desarrollo de \textbf{pipelines de datos} para plataformas de facturación electrónica, compras y ventas.
  \item Implementación de procesos \textbf{ETL/ELT} para ingestión, transformación y carga de datos desde SUNAT, pasarelas de pago y sistemas externos.
  \item Despliegue de soluciones de datos en \textbf{AWS} utilizando contenedores \textbf{Docker}.
  \item Optimización de consultas \textbf{SQL} y modelamiento de datos para mejorar los tiempos de respuesta en la capa de explotación.
  \item Automatización de controles de \textbf{calidad de datos} y validaciones de integridad.
  \item Monitoreo del consumo de recursos en la nube y optimización de costos.
\end{itemize}

\textbf{Tecnologías:} Python, SQL, AWS, Docker, CI/CD, Data Modeling

\vspace{10pt}

\textbf{Software Engineer III – Encora (Cliente: Profuturo AFP)} \hfill \textit{Nov 2020 – Abr 2025}
\begin{itemize}[leftmargin=*]
  \item Desarrollo y mantenimiento de procesos de integración de datos entre sistemas core financieros.
  \item Implementación y optimización de consultas \textbf{SQL} para reporting y explotación de información.
  \item Soporte a procesos batch y flujos de datos productivos.
  \item Participación en pases a producción y control de despliegues.
  \item Análisis de requerimientos funcionales y apoyo en el \textbf{modelamiento de datos}.
\end{itemize}

\textbf{Tecnologías:} SQL Server, PostgreSQL, Python, Azure, CI/CD

\vspace{10pt}

\textbf{Analista de Procesos / Software Engineer – Prosegur} \hfill \textit{Abr 2018 – Mar 2020}
\begin{itemize}[leftmargin=*]
  \item Desarrollo de soluciones de análisis y explotación de datos operativos.
  \item Implementación de \textbf{Web Audit}, plataforma para auditoría de abastecimiento de ATMs mediante análisis de información multimedia.
  \item Desarrollo de \textbf{Web PER}, sistema de monitoreo de rutas con GPS, mapas interactivos y alertas automáticas.
\end{itemize}

\textbf{Tecnologías:} SQL Server, MySQL, PHP Laravel

\vspace{10pt}

\textbf{Desarrollador Web \& Full Stack .NET – Medios Industriales} \hfill \textit{Abr 2014 – Mar 2018}
\begin{itemize}[leftmargin=*]
  \item Desarrollo del sistema \textbf{Control de Tiempos Operativos (CTO)} para análisis de jornadas laborales.
  \item Implementación de la plataforma \textbf{Balizas}, sistema de monitoreo GPS en tiempo real.
  \item Optimización de procesos operativos mediante análisis de datos.
\end{itemize}

\vspace{6pt}

% ================= DATA ENGINEERING =================
\section*{Data Engineering}

\begin{itemize}[leftmargin=*]
  \item Desarrollo y mantenimiento de \textbf{pipelines de datos} end-to-end.
  \item Procesos \textbf{ETL/ELT} para ingestión, transformación y carga de datos.
  \item Diseño y despliegue de soluciones en \textbf{Data Lake}.
  \item Optimización de performance en consultas y procesos de datos.
  \item Automatización de controles de \textbf{calidad de datos}.
  \item Uso de \textbf{Apache Spark y Databricks} a nivel intermedio.
  \item Gestión de pases a producción y monitoreo de procesos.
\end{itemize}

\vspace{6pt}

% ================= TECNOLOGIAS =================
\section*{Tecnologías}

\textbf{Data:} SQL, Python, ETL/ELT, Data Modeling\\
\textbf{Big Data:} Databricks, Apache Spark\\
\textbf{Cloud:} Azure (Data Lake, Storage, Azure ML), AWS\\
\textbf{CI/CD:} GitHub Actions, Azure DevOps\\
\textbf{Contenedores:} Docker\\
\textbf{Observabilidad:} Application Insights, Log Analytics\\
\textbf{Otros:} Git, Arquitectura Cloud, Optimización de Costos

\vspace{6pt}

% ================= EDUCACION =================
\section*{Educación}

\textbf{Universidad Nacional de Ingeniería (UNI)}\\
Ingeniero Industrial \hfill \textit{2004 – 2012}\\
CIP: 301084

\vspace{4pt}

\textbf{Universidad Privada del Norte}\\
Diplomado en Gerencia de TI con mención en Disrupción Tecnológica\\
\textit{Ene 2025 – Sep 2025}

\vspace{6pt}

% ================= CURSOS =================

\vspace{4pt}
\section*{Especialización Técnica y Cursos}
\begin{itemize}
   \item \textbf{Advanced Data Engineer (Diplomado)} – DMC Instituto\\
  \textit{Septiembre 2025 – Actual | 144 horas | Certificado en proceso}

  \item \textbf{Python Data Science} – DSRP Certification Program\\
  \textit{Octubre 2025 | 20 horas}

  \item \textbf{Gen AI Training Path} – Coursera Technical Track\\
  \textit{Marzo 2025 | 60 horas | Certificación Credly}

  \item \textbf{Ingeniero DevOps – Azure Fundamentals AZ-900}\\
  Udemy / Microsoft Azure Certification\\
  \textit{Noviembre 2024 | 9 horas}

  \item \textbf{Arquitectura Hexagonal C\#.NET} – Technical Track\\
  \textit{Octubre 2025 | 8.5 horas}

  \item \textbf{Experto en Business Intelligence}\\
  Sistemas UNI – Módulo de Base de Datos\\
  \textit{Febrero 2016 | 48 horas}

  \item \textbf{Administración de SQL Server}\\
  Sistemas UNI – Módulo de Base de Datos\\
  \textit{Junio 2015 | 96 horas}

  \item \textbf{Visual Basic .NET}\\
  Institución de Formación Técnica\\
  \textit{Septiembre 2011 | 48 horas}
\end{itemize}

% ================= IDIOMAS =================
\section*{Idiomas}

Español: Nativo\\
Inglés: Intermedio

\end{document}
